\chapter{Trabalhos Relacionados}
\label{cap:trabalhos-relacionados}

A aplicação de técnicas de aprendizado profundo à classificação automatizada do câncer de próstata tem sido amplamente investigada na literatura, com avanços expressivos, particularmente no uso de imagens de lâminas histopatológicas completas, estratégias baseadas na extração de \textit{patches} e mecanismos de agregação em múltiplas escalas. Esta seção apresenta uma revisão dos principais trabalhos relacionados, organizados de acordo com suas abordagens metodológicas.

\section{Abordagens Baseadas em WSI}

\citeonline{Xiang2023Prostate} propuseram o framework \textit{GCN-MIL} para diagnóstico e graduação do câncer de próstata em WSIs, baseado em aprendizado fracamente supervisionado com rótulos ao nível da lâmina. O método é estruturado em três etapas principais: (i) extração de \textit{patches} a partir das WSIs e obtenção de representações por meio de uma \textit{ResNet50} pré-treinada via aprendizado auto-supervisionado; (ii) modelagem das relações espaciais entre \textit{patches} utilizando uma Rede Neural de Grafos (GCN), na qual cada instância é representada como um nó conectado a seus vizinhos; e (iii) agregação global das representações por meio de um mecanismo de atenção no contexto de \textit{Multiple Instance Learning} (MIL), permitindo a predição do grau ISUP ao nível da lâmina. Adicionalmente, os autores incorporam uma estratégia de \textit{Robust Training}, que realiza a filtragem iterativa de amostras de alta incerteza para mitigar o impacto de rótulos ruidosos. O modelo foi treinado no conjunto de desafio PANDA, com mais de 10 mil lâminas, e avaliado em cenários internos e externos multicêntricos, alcançando desempenho competitivo, com coeficiente Kappa quadrático de até 0,931 no conjunto interno.

Apesar dos resultados expressivos, a abordagem apresenta limitações relevantes. O método possui elevada complexidade computacional, envolvendo etapas custosas, como a extração de milhares de \textit{patches} por WSI, a construção de grafos baseados em relações espaciais e o processamento por GCN, o que implica alto consumo de memória e tempo de processamento, dificultando sua aplicação em ambientes com recursos limitados. Outro ponto crítico refere-se à ausência de uma discussão explícita sobre a preservação da natureza ordinal dos graus ISUP no processo de treinamento. Embora o problema seja tratado como uma tarefa de classificação multiclasse, o trabalho não deixa claro se adotam-se mecanismos específicos para incorporar a estrutura hierárquica inerente à progressão da doença, o que pode limitar a capacidade do modelo de diferenciar adequadamente erros entre classes adjacentes e distantes.

\citeonline{pyramidalcnn} propuseram um sistema automatizado para graduação do câncer de próstata em WSIs de biópsias utilizando uma arquitetura Pyramidal CNN composta por três redes convolucionais operando em diferentes escalas de \textit{patch}, permitindo capturar características morfológicas em múltiplos níveis de contexto espacial. O pipeline de processamento incluiu equalização de histograma para normalização de intensidade, extração de patches a partir das WSIs e seleção com base na proporção mínima de tecido presente, seguida de classificação \textit{patch-wise} dos padrões de Gleason (ISUP 1-ISPU 5) e agregação das probabilidades para inferência do \textit{Gleason Score} e dos \textit{Grade Groups}. O conjunto experimental foi composto por 896 lâminas inteiras de biópsias prostáticas (750 para treinamento, 96 para validação e 50 para teste), fornecidas pelo \textit{Radboud University Medical Center} e analisadas no \textit{University of Louisville Hospital}, com rotulagem semi-automática realizada por patologistas especialistas. Os experimentos geraram milhões de \textit{patches} multiescala para treinamento das redes, e os resultados indicaram acurácia de aproximadamente 0,77 na classificação dos padrões, além de desempenho na identificação dos grupos prognósticos com recall entre 50\% e 75\% e precisão entre 50\% e 69\%, comparável ao desempenho humano em cenários similares. Contudo, o método apresenta limitações relacionadas ao uso de rotulação semiautomática, à ausência de validação externa independente, ao desempenho moderado em classes de maior complexidade e à dependência da estratégia de seleção de \textit{patches}, fatores que restringem a avaliação de generalização clínica e de robustez em ambientes multi-institucionais.

\citeonline{bhattacharyya2024efficient} propôs dois métodos complementares para a graduação ISUP de câncer de próstata a partir de WSIs. O primeiro consiste em um ensemble hierárquico de CNNs (\textit{Multi-CNN Ensemble}), no qual um modelo global $M_1$, treinado via MIL (\textit{Multiple Instance Learning}) com regressão ordinal, realiza uma classificação inicial da WSI. A partir da análise dos pares de graus com o maior índice de confusão e os autores treinaram modelos especializados nos subconjuntos de graus confusos, com diferentes tamanhos de \textit{patch} ($192\times192$, $224\times224$ e $256\times256$). A predição final é obtida por meio de uma estratégia de decisão em cascata entre os modelos globais e especializados. O segundo método (\textit{Bag-of-CNN Features}) utiliza as representações extraídas por modelos CNN treinados para gerar, por meio de \textit{k}-means, um descritor global para WSI, que alimenta classificadores baseados em árvores de decisão com \textit{gradient boosting} (XGBoost e LGBM). Ambos os métodos adotam o EfficientNet-B1 como \textit{backbone} e são avaliados no dataset PANDA, no qual o \textit{Multi-CNN Ensemble} alcançou kappa de $0{,}881$.

\citeonline{kabir2025automating} propuseram um framework de três estágios que integra classificação e segmentação em nível de patch para a predição do grau ISUP a partir de WSIs. Após uma etapa de limpeza de dados, patches de dois tamanhos distintos (500x500 e 1000x1000 \textit{pixels}) foram extraídos para alimentar as redes de classificação e de segmentação. No primeiro estágio, modelos de classificação baseados em EfficientNet-b0, DenseNet121, Inception-v3 e Vision Transformer foram treinados para classificar os \textit{patches} em quatro classes: benigno, Gleason 3, Gleason 4 e Gleason 5. No segundo estágio, uma arquitetura de segmentação semântica DeepLabV3+, com codificador EfficientNet-B0 e camadas convolucionais substituídas por redes neurais operacionais auto-organizáveis (Self-ONN), foi proposta para a delimitação \textit{pixel} a \textit{pixel} de cada padrão histológico. No terceiro estágio, as proporções de cada padrão de Gleason nas WSIs foram concatenadas em vetores de características para o treinamento de classificadores de aprendizado de máquina. O classificador Random Forest, treinado com a combinação de características provenientes de ambas as abordagens e de ambos os tamanhos de patch, obteve o melhor kappa quadrático de 0,921 no \textit{dataset} PANDA. Apesar dos resultados promissores, o estudo apresenta uma limitação metodológica importante ao descartar as imagens do Instituto \textit{Karolinska}. A utilização exclusiva de dados de uma única instituição (\textit{Radboud}) impede a validação externa do modelo , restringindo a sua capacidade de generalização para dados inéditos e, consequentemente, limitando a sua viabilidade para adoção na prática clínica.

\citeonline{bhattacharyya2026tsor} propuseram o TSOR (\textit{Task-Specific Ordinal Regression}), um \textit{framework} composto por três módulos integrados. No primeiro, um modelo híbrido de MIL gera pseudo-rótulos em nível de \textit{patch}, produzindo um conjunto de dados relativamente balanceado entre os graus de Gleason no pré-treinamento. No segundo, esse conjunto é utilizado em um \textit{framework} de aprendizado auto-supervisionado específico para a tarefa, com um termo de perda adicional baseado em augmentação de coloração, permitindo o aprendizado de características morfológicas invariantes à variação de \textit{stain} entre os centros. No terceiro módulo, a rede pré-treinada é ajustada para a graduação ISUP por meio de regressão ordinal, na qual os rótulos são codificados de modo que a distância entre as classes reflita sua ordenação clínica. Avaliado nos conjuntos PANDA e SICAPv2, o método atingiu kappa de $0{,}884$ e $0{,}874$, respectivamente, superando abordagens SOTA baseadas em MIL e modelos de fundação histopatológicos.

\section{Abordagens Baseadas em Patches}

No trabalho de \citeonline{kosoko2024xai_prostate}, os autores avaliaram diferentes arquiteturas de redes neurais convolucionais, combinadas com técnicas de \textit{Explainable Artificial Intelligence} (XAI), em especial \textit{Grad-CAM}, para a interpretação das decisões do modelo. O conjunto de dados utilizado foi o SICAPv2. Entre as arquiteturas avaliadas, o modelo VGG19 apresentou o melhor desempenho, alcançando AUC-ROC de 0{,}85 e valores médios próximos de 0{,}75 para acurácia, precisão e \textit{recall}, enquanto uma CNN simples apresentou desempenho significativamente inferior (acurácia $\approx$ 0{,}41 e F1 $\approx$ 0{,}24), evidenciando a vantagem de arquiteturas profundas pré-treinadas na captura de padrões morfológicos complexos associados aos graus de Gleason. O estudo apresenta limitações relacionadas à ausência de validação clínica externa, à dependência de um único conjunto de dados e à exploração restrita de técnicas de XAI, o que limita a avaliação de generalização em cenários interinstitucionais.


\citeonline{alici2024efficient} propuseram um modelo de aprendizado profundo para diagnóstico e graduação do câncer de próstata utilizando imagens histopatológicas do dataset público DiagSet. O processamento foi realizado em \textit{patches} extraídos de lâminas digitalizadas coradas em H\&E. A abordagem baseou-se na arquitetura EfficientNet-B4 e em uma variante denominada Eff4-Attn, que incorpora um mecanismo de atenção de canais para aprimorar a extração de características morfológicas. O modelo foi avaliado em diferentes níveis de granularidade: uma classificação em quatro classes (tecido normal, estroma, alterações benignas e câncer) e outra em oito classes (tecido normal, estroma, artefatos, inflamação, hiperplasia benigna, Gleason 3, Gleason 4 e Gleason 5). Os resultados alcançaram acurácia de 96,18\% na tarefa binária, 94,86\% na classificação em quatro classes e 93,32\% na classificação em oito classes, demonstrando desempenho competitivo aliado à eficiência computacional. Contudo, o estudo apresenta limitações relacionadas ao uso exclusivo de \textit{patches} em vez de inferência em nível de lâmina completa, à dependência de um único conjunto de dados público e ao desbalanceamento entre as classes, fatores que restringem a avaliação da generalização clínica e da robustez do método em cenários multi-institucionais.

No estudo de \citeonline{ikromjanov2021vit_prostate}, a análise de imagens histopatológicas de próstata foi conduzida por meio de um modelo baseado em \textit{Vision Transformers} (ViT), com o objetivo de classificar padrões teciduais associados ao sistema de graduação de Gleason a partir de WSI. O conjunto de dados utilizado deriva do desafio público PANDA, sendo empregado um subconjunto com mais de 5.000 imagens na experimentação. Após a detecção das regiões de interesse, as lâminas foram particionadas em aproximadamente 304 mil \textit{patches} de dimensões 256×256 pixels, rotuladas em cinco classes (estroma, benigno, Gleason 3, Gleason 4 e Gleason 5), com cerca de 240 mil amostras destinadas ao treinamento, 60 mil à validação e 4 mil ao teste. O modelo ViT foi utilizado para classificação em nível de \textit{patch} e para inferência em nível de lâmina, com predições agregadas. Os resultados experimentais demonstraram desempenho competitivo na identificação de padrões histológicos, alcançando acurácia superior a 85\% na classificação dos \textit{patches}, evidenciando a capacidade do mecanismo de atenção de capturar dependências contextuais globais nas imagens. 

\section{Limitações da Literatura e Lacunas de Pesquisa}

Apesar dos avanços no uso de aprendizado profundo na análise de imagens histopatológicas de próstata, persistem limitações relevantes. Muitos trabalhos tratam o problema como uma classificação nominal, ignorando a natureza ordinal dos graus ISUP, o que penaliza erros de forma uniforme e desalinha o modelo em relação à relevância clínica. Além disso, há forte dependência de arquiteturas convolucionais e de estratégias baseadas em \textit{patches}, sem tratamento robusto de ruído e incerteza, o que é agravado por artefatos, inconsistências nas anotações e variabilidade entre instituições em \textit{datasets} como o PANDA.

Outra limitação é o uso predominante de classificadores individuais, com pouca exploração de \textit{ensembles}, além da dificuldade em capturar relações hierárquicas entre graus tumorais. Para enfrentar essas lacunas, este trabalho propõe um pipeline que combina a filtragem por entropia, a modelagem ordinal com \textit{Focal Loss} e um \textit{ensemble} de variantes do EfficientNet, visando maior robustez, estabilidade e coerência clínica nas predições.		